\section*{Chapitre \ref{ch:g10:layered-ocean} --- Un océan en couches}

\begin{solution}{exo:g10:layered-ocean:1}
La \omterm{def:g10:layered-ocean:mixed-layer}{couche de mélange} est la couche de surface de l'océan dans laquelle la
température et la salinité varient à peine avec la profondeur~; la
\omterm{def:g10:layered-ocean:thermocline}{thermocline} est la couche située en dessous, dans laquelle la température
chute vite avec la profondeur. Le vent et les vagues brassent la \omterm{def:g10:layered-ocean:mixed-layer}{couche de
mélange} (aidés en hiver par le refroidissement de la surface).
\end{solution}

\begin{solution}{exo:g10:layered-ocean:2}
Une \omterm{def:g10:layered-ocean:upwelling}{remontée d'eau} est la montée d'eau froide venue de sous la couche de
surface, qui remplace l'eau de surface repoussée par le vent. L'eau qui
remonte apporte des sels minéraux nutritifs dans la couche éclairée, où le
phytoplancton les utilise pour croître~; le phytoplancton nourrit le reste
de la chaîne alimentaire.
\end{solution}

\begin{solution}{exo:g10:layered-ocean:3}
Pacifique tropical~: environ \qty{29}{\degreeCelsius} en surface et
\qty{4.5}{\degreeCelsius} à \qty{1000}{m}. Océan Arctique~: environ
\qty{-1.4}{\degreeCelsius} en surface et près de \qty{0}{\degreeCelsius} à
\qty{1000}{m}.
\end{solution}

\begin{solution}{exo:g10:layered-ocean:4}
À \qty{0}{\degreeCelsius} (\qty{1028.1}{kg/m^3} contre \qty{1023.3}{kg/m^3}
à \qty{25}{\degreeCelsius}). L'eau de mer~: à \qty{4}{\degreeCelsius}, l'eau
pure a une masse volumique de \qty{1000.0}{kg/m^3}, l'eau de mer de
salinité 35 environ \qty{1027.8}{kg/m^3} (lu sur le graphique).
\end{solution}

\begin{solution}{exo:g10:layered-ocean:5}
Le flotteur dérive à \qty{1000}{m}, plonge à \qty{2000}{m}, puis remonte à
la surface en mesurant au passage la température, la salinité (déduite de
la conductivité) et la pression~; en surface, il transmet son profil par
satellite et replonge. Les flotteurs dérivent avec les courants dans tout
l'océan et fonctionnent pendant des années, en toutes saisons, loin des
routes maritimes.
\end{solution}

\begin{solution}{exo:g10:layered-ocean:6}
$361.84\times10^6 \times 3.682 \approx \qty{1.33e9}{km^3}$, environ
\num{1.33}~milliard de kilomètres cubes. $100/3682 \approx 0.027$~: moins
de \qty{3}{\percent} de la profondeur moyenne.
\end{solution}

\begin{solution}{exo:g10:layered-ocean:7}
$22.7 - 0.5 = \qty{22.2}{\degreeCelsius}$, atteint à \qty{15}{m}~: la \omterm{def:g10:layered-ocean:mixed-layer}{couche
de mélange} a environ \qty{15}{m} de profondeur. Entre \qty{20}{m} et
\qty{50}{m}~:
$(21.8 - 17.3)/30 = \qty{0.15}{\degreeCelsius}$ par mètre.
\end{solution}

\begin{solution}{exo:g10:layered-ocean:8}
Le soleil ne chauffe que la couche de surface, et le réchauffement et le
refroidissement saisonniers ne se répartissent que dans la \omterm{def:g10:layered-ocean:mixed-layer}{couche de
mélange}~: sous la \omterm{def:g10:layered-ocean:thermocline}{thermocline}, l'eau en est coupée. En hiver, la surface se
refroidit, devient plus dense et plonge dans l'eau sous-jacente, et les
tempêtes brassent la mer~: le mélange descend profond. En été, le soleil
réchauffe une couche de surface légère qui flotte sur l'eau plus froide~;
les vents, plus doux, ne peuvent l'enfoncer, et la \omterm{def:g10:layered-ocean:mixed-layer}{couche de mélange} reste
mince.
\end{solution}

\begin{solution}{exo:g10:layered-ocean:9}
$0.78/0.18 \approx 4.3$ degrés vers \qty{10}{\degreeCelsius}~;
$0.78/0.06 = 13$ degrés vers \qty{0}{\degreeCelsius}. Dans les mers
polaires froides, de petites variations de salinité comptent bien plus pour
la masse volumique que les variations de température~: c'est la salinité
qui y commande l'étagement.
\end{solution}

\begin{solution}{exo:g10:layered-ocean:10}
La température est la même~: l'eau profonde est plus dense parce qu'elle
est plus salée (33.3 contre 28.1). La couche qui les sépare, où la salinité
change vite avec la profondeur, est une \omterm{def:g10:layered-ocean:thermocline}{halocline}~; ici, c'est aussi la
\omterm{def:g10:layered-ocean:thermocline}{pycnocline}.
\end{solution}

\begin{solution}{exo:g10:layered-ocean:11}
Poussée vers le bas, la parcelle garde sa température et sa salinité, donc
sa masse volumique. Entourée d'eau plus dense, elle est plus légère que l'eau
qu'elle déplace~: la poussée d'Archimède la ramène à son propre niveau. Une
parcelle poussée vers le haut serait plus dense que son nouvel environnement
et redescendrait. Les couches sont stables.
\end{solution}

\begin{solution}{exo:g10:layered-ocean:12}
Masse~: $0.001 \times 1028.1 = \qty{1.0281}{kg}$. Nouveau volume~: $1.0281/1023.3
\approx \qty{0.0010047}{m^3}$, environ \qty{1004.7}{mL}~: il s'est dilaté
d'environ \qty{4.7}{mL}, soit \qty{0.47}{\percent}. La couche de
\qty{700}{m} s'épaissirait de $700 \times 0.0047 \approx \qty{3.3}{m}$.
\end{solution}

\begin{solution}{exo:g10:layered-ocean:13}
Dans l'Atlantique Nord en février. Les \qty{150}{m} supérieurs y ont presque
la même température, donc presque la même masse volumique~: les brasser
coûte peu. Dans le Pacifique tropical, la masse volumique augmente de
$1026.0 - 1021.6 =
\qty{4.4}{kg/m^3}$ entre \qty{50}{m} et \qty{200}{m}~: une \omterm{def:g10:layered-ocean:thermocline}{pycnocline}
marquée, et la tempête devrait enfoncer de l'eau légère dans une eau bien
plus dense.
\end{solution}

\begin{solution}{exo:g10:layered-ocean:14}
Chlorophylle abondante loin des côtes~: la bande qui longe l'équateur dans
le Pacifique (et dans l'Atlantique), où les alizés entretiennent une
\omterm{def:g10:layered-ocean:upwelling}{remontée d'eau}~; et les hautes latitudes, comme l'Atlantique Nord, où la
\omterm{def:g10:layered-ocean:mixed-layer}{couche de mélange} d'hiver descend profond et fait remonter les nutriments.
Chlorophylle très faible~: le centre des gyres subtropicaux, fortement
\omterm{def:g10:layered-ocean:stratification}{stratifiés}, où rien ne fait remonter l'eau profonde et où la couche éclairée
s'est épuisée en nutriments.
\end{solution}

\begin{solution}{exo:g10:layered-ocean:15}
Plongée~: $1000/(6 \times 3600) \approx \qty{0.046}{m/s}$, environ
\qty{4.6}{cm/s}. Remontée~: $2000/(10 \times 3600) \approx \qty{0.056}{m/s}$,
environ \qty{5.6}{cm/s}. $36\,000/2 = 18\,000$ mesures, espacées en moyenne
de $2000/18\,000 \approx \qty{0.11}{m}$.
\end{solution}

\begin{solution}{pb:g10:layered-ocean:1}
\textbf{1.} Le gyre, de $24.0 - 19.0 = \qty{5.0}{\degreeCelsius}$.

\textbf{2.} Le gyre (\qty{20.5}{\degree S} contre \qty{13.5}{\degree S}).
Plus éloigné de l'équateur, il devrait être le plus frais~: les seules
latitudes ne peuvent expliquer qu'il soit plus chaud de
\qty{5}{\degreeCelsius}.

\textbf{3.} $24.0 - 0.5 = \qty{23.5}{\degreeCelsius}$, atteint entre
\qty{50}{m} (23.6) et \qty{100}{m} (22.6), vers $50 + 50 \times
0.1/1.0 = \qty{55}{m}$.

\textbf{4.} Non~: la température baisse régulièrement dès la surface, de
\qty{1.6}{\degreeCelsius} sur les \qty{25}{m} supérieurs et de
\qty{5.3}{\degreeCelsius} sur les \qty{100}{m} supérieurs~; la \omterm{def:g10:layered-ocean:thermocline}{thermocline}
commence à la surface.

\textbf{5.} Entre \qty{200}{m} et \qty{300}{m}~: $20.3 - 14.4 =
\qty{5.9}{\degreeCelsius}$ par centaine de mètres.

\textbf{6.} Au large du Pérou, entre \qty{50}{m} et \qty{100}{m}~: environ
\qty{70}{m} par interpolation (le profil complet donne \qty{65}{m}). Dans le
gyre, entre \qty{200}{m} et \qty{300}{m}~: environ $200 + 100 \times 5.3/5.9
\approx \qty{290}{m}$.

\textbf{7.} Oui, à tous les niveaux des deux colonnes~: l'eau plus légère
repose sur l'eau plus dense, et les deux colonnes sont stables.

\textbf{8.} Une eau plus salée devrait être plus dense~; pourtant l'eau de
surface du gyre est plus légère~: \qty{1024.71}{kg/m^3} contre
\qty{1025.08}{kg/m^3}.

\textbf{9.} L'eau du gyre est plus salée d'environ 1.3 unité, ce qui ajoute
environ $1.3 \times 0.8 \approx \qty{1.0}{kg/m^3}$~; elle est plus chaude de
\qty{5}{\degreeCelsius}, ce qui retire environ $5 \times 0.27 \approx \qty{1.4}{kg/m^3}$ à ces
températures. La température l'emporte.

\textbf{10.} $1026.44 - 1025.08 = \qty{1.36}{kg/m^3}$.

\textbf{11.} $1025.30 - 1024.71 = \qty{0.59}{kg/m^3}$.

\textbf{12.} De la surface à \qty{500}{m}, la masse volumique augmente de
$1026.92 - 1024.71 = \qty{2.21}{kg/m^3}$, dont seulement \qty{0.59}{kg/m^3}
au-dessus de \qty{200}{m}~: l'essentiel, \qty{1.62}{kg/m^3}, se trouve sous
\qty{200}{m}, dans la \omterm{def:g10:layered-ocean:thermocline}{thermocline} entre \qty{200}{m} et \qty{500}{m}.

\textbf{13.} Au large du Pérou, la masse volumique augmente de $1026.97 - 1025.08 =
\qty{1.89}{kg/m^3}$ de la surface à \qty{500}{m}, dont
\qty{1.36}{kg/m^3} dans les \qty{200}{m} supérieurs~: l'essentiel se trouve
près de la surface.

\textbf{14.} À \qty{500}{m}, les deux sites sont très semblables
(\qty{7.9}{\degreeCelsius} et \qty{7.1}{\degreeCelsius},
\qty{1026.97}{kg/m^3} et \qty{1026.92}{kg/m^3})~: la même eau froide
s'étend sous les deux~; la différence tient à la distance à laquelle elle
s'approche de la surface.

\textbf{15.} Le vent qui souffle le long du rivage vers l'équateur repousse
l'eau de surface vers le large~; de l'eau venue d'en dessous remonte près
de la côte pour prendre sa place. L'eau froide de la \omterm{def:g10:layered-ocean:thermocline}{thermocline} est
aspirée vers le haut, si bien que la \omterm{def:g10:layered-ocean:thermocline}{thermocline} remonte en pente vers la
côte et atteint presque la surface.

\textbf{16.} Au large du Pérou~: de l'eau à \qty{13}{\degreeCelsius} à
\qty{15}{\degreeCelsius}, venue de sous la \omterm{def:g10:layered-ocean:thermocline}{thermocline}, se trouve dans les
\qty{100}{m} supérieurs, à portée du brassage par le vent et de la lumière.
Dans le gyre, elle se trouve sous \qty{250}{m}, sous une couche chaude et
légère que le vent ne peut brasser aussi profond.

\textbf{17.} La \omterm{def:g10:layered-ocean:upwelling}{remontée d'eau} cesserait~: l'eau de surface se réchaufferait
au soleil, une couche chaude et légère se formerait en haut, et l'apport de
nutriments à la couche éclairée s'amenuiserait.

\textbf{18.} La côte du Pérou est bleu clair à jaune (beaucoup de
chlorophylle)~; le milieu du gyre du Pacifique Sud compte parmi les régions
les plus foncées de la carte (très peu).

\textbf{19.} Le phytoplancton a besoin de nutriments autant que de lumière.
Dans le gyre, la forte \omterm{def:g10:layered-ocean:stratification}{stratification} maintient l'eau profonde, riche en
nutriments, bien en dessous de la couche éclairée, et les nutriments de la
couche de surface sont épuisés depuis longtemps.

\textbf{20.} $1.36/0.59 \approx 2.3$~: \textbf{l'augmentation de masse
volumique sur les \qty{200}{m} supérieurs est environ 2.3 fois plus grande
au large du Pérou que dans le gyre}. La \omterm{def:g10:layered-ocean:thermocline}{pycnocline}, avec l'eau froide et
riche en nutriments qui se trouve sous elle, a été soulevée dans les
\qty{200}{m} supérieurs au large du Pérou, tandis que dans le gyre elle se
trouve plus profond, sous \qty{200}{m}.
\end{solution}
